\documentclass[11pt,a4paper]{article}

% ---------------------------------------------------------------
%  TD1 -- Structures arborescentes -- Exercice 2
%  Compilation : pdflatex TD1_Structures_Arborescentes_Ex2.tex
% ---------------------------------------------------------------

\usepackage[utf8]{inputenc}
\usepackage[T1]{fontenc}
\usepackage[french]{babel}
\usepackage{lmodern}
\usepackage[margin=2.3cm]{geometry}
\usepackage{amsmath,amssymb}
\usepackage{array}
\usepackage{xcolor}
\usepackage{listings}
\usepackage{enumitem}

% ---------- Mise en forme du code ----------
\definecolor{codebg}{RGB}{248,248,245}
\definecolor{codekw}{RGB}{0,70,160}
\definecolor{codestr}{RGB}{160,40,40}
\definecolor{codenum}{RGB}{130,130,130}

\lstdefinestyle{td}{
  language=Python,
  basicstyle=\ttfamily\small,
  keywordstyle=\color{codekw}\bfseries,
  stringstyle=\color{codestr},
  numberstyle=\tiny\color{codenum},
  numbers=left,
  numbersep=8pt,
  backgroundcolor=\color{codebg},
  frame=single,
  rulecolor=\color{black!20},
  framesep=6pt,
  xleftmargin=18pt,
  breaklines=true,
  showstringspaces=false,
  tabsize=4,
  columns=fullflexible,
  keepspaces=true,
  morekeywords={int,bool,True,False,None},
  literate={é}{{\'e}}1 {è}{{\`e}}1 {à}{{\`a}}1
}
\lstset{style=td}

% ---------- Macros pour les monômes ----------
% Monôme (coeff, degré)
\newcommand{\mono}[2]{%
  $\left\{\begin{array}{@{}l@{\;\;}l@{}}
    \text{coeff} & #1\\
    \text{degré} & #2
  \end{array}\right.$}

% Monôme avec indicateur de sentinelle (coeff, estSent, degré)
\newcommand{\monoS}[3]{%
  $\left\{\begin{array}{@{}l@{\;\;}l@{}}
    \text{coeff}   & #1\\
    \text{estSent} & #2\\
    \text{degré}   & #3
  \end{array}\right.$}

% Flèche « ondulée » utilisée dans les notes
\newcommand{\vers}{\;\rightsquigarrow\;}

% Titres des questions
\newcommand{\question}[1]{\par\medskip\noindent\textbf{#1}\quad}

\renewcommand{\arraystretch}{1.25}
\setlength{\parindent}{0pt}

\title{\textbf{TD1 -- Structures arborescentes}\\[4pt]\large Exercice 2}
\author{}
\date{}

\begin{document}
\maketitle
\vspace{-4em}

% ===============================================================
\section*{Exercice 2}

% ---------------------------------------------------------------
\question{Q1 /}
Un polynôme est un tableau de monômes :
\[
  x^4 + 3x^{100} \vers
  \begin{array}{|c|c|}
    \multicolumn{1}{c}{0} & \multicolumn{1}{c}{1}\\ \hline
    x^4 & 3x^{100} \\ \hline
  \end{array}
\]

\begin{lstlisting}
def polynome(p: Array[monomes]) -> None:
    int nb_termes = len(p)

    for i in range(0, nb_termes):

        p[i].coeff = p[i].degre * p[i].coeff
        if p[i].degre != 0:
            p[i].degre = p[i].degre - 1
\end{lstlisting}

% ---------------------------------------------------------------
\question{Q2 /}
$x \mapsto x^4 + 3x^{100}$

\subsection*{Représentation 1}

\begin{center}
monomes $\vers$\quad
\begin{tabular}{|c|c|}
  \multicolumn{1}{c}{0} & \multicolumn{1}{c}{1}\\ \hline
  \mono{3}{100} & \mono{1}{4} \\ \hline
\end{tabular}
\end{center}

\subsection*{Représentation 2 \normalfont(1 sentinelle)}

\begin{center}
monomes $\vers$\quad
\begin{tabular}{|c|c|c|}
  \multicolumn{1}{c}{0} & \multicolumn{1}{c}{1} & \multicolumn{1}{c}{2}\\ \hline
  \monoS{0}{1}{0} & \monoS{3}{0}{100} & \monoS{1}{0}{4} \\ \hline
\end{tabular}
\end{center}

\subsection*{Représentation 3 \normalfont(2 sentinelles)}

\begin{center}
monomes $\vers$\quad
\begin{tabular}{|c|c|c|c|}
  \multicolumn{1}{c}{0} & \multicolumn{1}{c}{1} &
  \multicolumn{1}{c}{2} & \multicolumn{1}{c}{3}\\ \hline
  \monoS{0}{1}{0} & \monoS{3}{0}{100} & \monoS{1}{0}{4} & \monoS{0}{1}{0} \\ \hline
\end{tabular}
\end{center}

\subsection*{Représentation 4}
% L'indice de la case correspond au degré : toutes les cases sont
% initialisées avec le monôme nul, sauf celles des degrés 4 et 100.

\begin{center}
monomes $\vers$\quad
\begin{tabular}{|c|c|c|c|c|c|}
  \multicolumn{1}{c}{0} & \multicolumn{1}{c}{1} & \multicolumn{1}{c}{}  &
  \multicolumn{1}{c}{4} & \multicolumn{1}{c}{}  & \multicolumn{1}{c}{100}\\ \hline
  \mono{0}{0} & \mono{0}{0} & $\cdots\cdots$ & \mono{1}{4} & $\cdots\cdots$ & \mono{3}{100} \\ \hline
\end{tabular}
\end{center}

% ---------------------------------------------------------------
\question{Q3 /}
Représentation canonique du polynôme
$x \mapsto 1\cdot x^{100} + 0\cdot x^{6} + x^{4} + 2\cdot x^{100}$ :

\begin{center}
monomes $\vers$\quad
\begin{tabular}{|c|c|}
  \multicolumn{1}{c}{0} & \multicolumn{1}{c}{1}\\ \hline
  \mono{1}{4} & \mono{3}{100} \\ \hline
\end{tabular}
\end{center}

% ---------------------------------------------------------------
\question{Q4 /}

\begin{itemize}[leftmargin=2em, itemsep=2pt]
  \item[] \textbf{Auxiliaire 1 :} vérifier le coeff non nul.
  \item[] \textbf{Auxiliaire 2 :} vérifier si le degré est déjà dans la table
          et additionner leurs coefficients.
  \item[] \textbf{Auxiliaire 3 :} tri des monômes par degré croissant.
\end{itemize}

% ---------------------------------------------------------------
\question{Q5 /}

\begin{lstlisting}
def sont_egaux(pol1: Array[monomes], pol2: Array[monomes]) -> bool:

    pol1 = canonique(pol1)
    pol2 = canonique(pol2)

    int t1 = size(pol1)
    int t2 = size(pol2)
    if t1 != t2:
        return False
    int res = True
    for i in range(0, t1):

        if (pol1[i].degre != pol2[i].degre) or (pol1[i].coeff != pol2[i].coeff):
            res = False

    return res
\end{lstlisting}

\end{document}